% !PRINTED 536
\weberpage{536}

Aus einer solchen Basisform erhält man alle ganzen Zahlen des Körpers
$\Omega$, wenn man für die Variablen ganze rationale Zahlen setzt.

Die Grundform $\tau$ ist die Wurzel einer irreduciblen Gleichung
$n^{\text{ten}}$ Grades
\[
F(t) = N(t - \tau) = 0,
\]
in der die Co\"efficienten der Potenzen von $t$ ganze rationale (und
homogene) Functionen der Variablen $t_1, t_2, \ldots, t_n$ sind.

\begin{center}
\S. 147.

\gesperrt{Die absoluten Normen der Functionale.}
\end{center}

Die Elemente $\alpha_1, \alpha_2, \ldots, \alpha_n$ einer Basis des
Functionals $\mu$ können als ganze Zahlen in $\Omega$ linear und
ganzzahlig ausgedrückt werden durch eine Basis $\omega_1, \omega_2,
\ldots, \omega_n$ von $\rho$ in der Form
\weq{1}{
(\alpha_1, \alpha_2, \ldots, \alpha_n) = A(\omega_1, \omega_2, \ldots,
\omega_n),
}
worin $A$ eine lineare Substitution mit ganzen rationalen
Co\"efficienten bedeutet. Einen Specialfall hiervon bieten die Formeln
(3) des vorigen Paragraphen.

Eine Basisform $\lambda$ von $\mu$ wollen wir so darstellen:
\weq{2}{
\lambda = \sum^{\nu} \alpha_\nu t_\nu,
}
und die Substitution (1) schreiben wir ausführlicher:
\weq{3}{
\alpha_s = \sum^{\nu} a_{s,\nu}\,\omega_\nu,
}
worin $t_\nu$ Variable, $a_{s,\nu}$ die Substitutionsco\"efficienten
sind, und der Summationsbuchstabe $\nu$ von 1 bis $n$ läuft.

Da die Producte $\alpha_s\,\omega_r$ alle durch $\mu$ theilbar sind, so
können sie nach der Bedeutung der Basis in der Form dargestellt werden
\weq{4}{
\alpha_\nu\,\omega_r = \sum^{s} g^{(s)}_{\nu,r}\,\alpha_s,
}
worin die $g^{(s)}_{\nu,r}$ ganze rationale Zahlen sind. Hieraus erhält
man dann nach (2)
\weq{5}{
\lambda\,\omega_r = \sum^{s} \alpha_s\,t_{s,r},
}
wenn
\weq{6}{
t_{s,r} = \sum^{\nu} g^{(s)}_{\nu,r}\,t_\nu
}
ganze rationale Linearformen sind.

Substituirt man in (5) wieder die Ausdrücke (3), so folgt
\weq{7}{
\lambda\,\omega_r = \sum^{\nu} \omega_\nu \sum^{s} a_{s,\nu}\,t_{s,r}.
}
